% Transcription of folder 147, batch 09 (archivists' pages 161-180).
% Pass: Opus 5.5 (claude-opus-5-5), 2026-10-09 — first pass, unchecked against the pages by a human.
%
% Author's pp. 77-86 (two leaves both numbered 81). Fast working hand; the
% formulas carry better than the connective prose. Notation fixed for the
% batch: D_{d_*}, D^*_{d_*} strata, V_{A,R} / V^*_{A,R} / V^{d'_0}_{d_*,d''_0}
% tubular-neighbourhood topoi, \check{N} for his normal bundle, \Pi_1
% fundamental groupoid, \mathbb{D}, \mathbb{D}^* germs of discs, \mathbb{P} point.
\documentclass[11pt,a4paper]{article}
\input{../preamble/grothendieck}

\folder{147}
\batch{9}
\pages{161}{180}
\foldertitle{Structure à l'infini des Mg,ν [pages 1 à 90, dont table des matières] : notes manuscrites (s.d.).}
\dating{s.d. — le groupe « Autour de La "Longue Marche" à travers la théorie de Galois » (141 à 149) est daté [à partir de 1978]-1983}
\watermark{Édition de démonstration}

\begin{document}

\page{161}\note{p.~77 de l'auteur (le chiffre est surchargé). La phrase qui ouvre la page continue un argument commencé avant ce lot.}
Dans un tel cas, les topos envisagés comme « pièces de construction » sont les $\mathcal{V}_{d_*,d'_*}$ et $\mathcal{V}^*_{d_*,d'_*}$. Du p\supplied{oin}t de vue notation, j'ai omis d'éliminer des redondances, telles que $D_{d_0,d_0}\simeq D_{d_0}$ et $D_{d_0,d_1,d_1}\simeq D_{d_0,d_1}$ (les op\supplied{érations} \add{simpl.} de dégénérescence sont des iso\,!) en exigeant $d_0>\dots>d_r$ dans la notation $D_{d_*}=D_{d_0,\dots,d_r}$ (ce qui indique qu'il \ill{} \ill{}\,: des op. simpliciales qui correspondent à des inclusions de simplexes types --- donc\,: des composés d'op. face), \struck{donc} et les redondances du type $D_{d,0}\simeq D_d$, plus gén. $D_{d_0,\dots,d_r,0}\simeq D_{d_0,\dots,d_r}$ (si $d_r>0$) en exigeant que si $r\geqslant 1$, alors $d_r>0$. \struck{\ill{}} \uncertain{Avec} cette notation, le seul $D_{d_*}$ où figure l'indice $0$ est $D_0$. Il faudrait donc compléter \uncertain{implicitement} les op. simpliciales entre les $D_{d_*}$ en y ajoutant les projections
\[ D_{d_*}\longrightarrow D_0=X \]
p.\,ex. en regardant les $D_{d_*}$ comme étant des espaces analytiques \emph{au-dessus de $X$} (les op. simpliciales respectant cette structure).
\marginal{Ces morphismes correspondraient formellement à l'inclusion du système \emph{vide} $\emptyset\subset\{d_0>\dots>d_r\}$.}

\page{162}\note{p.~78 de l'auteur.}
\marginal{Dans tous les cas, les topos \uncertain{composants} \struck{\ill{}} envisagés sont des cas particuliers des $\mathcal{V}_{d_*,d'_*}$ et $\mathcal{V}^*_{d_*,d'_*}$ (260).}
Je vais faire la liste complète des \add{typiques} \struck{\ill{}}, en topos, dans le cas le plus élémentaire de tous\,:
\[ (261)\qquad X=\mathbb{E}^{1\,\mathrm{an}},\qquad \Theta=\{0\} \]
on aura donc (en \uncertain{n'incluant} que les $D_{d_*}$ \uncertain{étant} non vides) les « drapeaux » $\{0\}$, $\{1\}$, $\{1,0\}$ et les inclusions de drapeaux non triviales $\{0\}\to\{0,1\}$, $\{1\}\to\{0,1\}$.

On aura \struck{$D_0=\mathcal{V}_0=\mathbb{C}$, $D_0^*=\mathbb{C}^*$ ($\mathcal{V}_0^*$), $D_1=D_1^*\simeq D_{1,0}=D_{1,0}^*=$ origine de $\mathbb{C}$, $\mathcal{V}_1=\mathcal{V}_{0,\mathbb{C}}$}\note{bloc biffé de traits obliques, en partie raturé au point d'être illisible\,; la lecture ne vaut que pour ce qui reste visible.}
\[ (262)\quad \left\{\begin{array}{l} D_0=\mathbb{C},\qquad \boxed{D_0^*=\mathbb{C}^*}\\[2pt] D_{1,0}\xrightarrow{\ \sim\ } D_1=\{\mathbb{P}\},\qquad D_{1,0}^*=D_{1,0}\\ \quad\big\downarrow\ \text{inclusion de $\{0\}$ dans $\mathbb{C}$}\\ D_0=\mathbb{C} \end{array}\right. \qquad D_1=\boxed{D_1^*=\{\mathbb{P}\}} \]
\note{à droite de $D_1^*=\{\mathbb{P}\}$, biffé\,: « orig\ill{} », « (origine de $\mathbb{C}$) ». Le numéro (262) est supposé\,: seul (263) se lit, dans la marge.}
\[ (263)\quad \left\{\begin{array}{ll} \boxed{\mathcal{V}_0^*\simeq\mathcal{V}_0\xrightarrow{\ \simeq\ } D_0^*=\mathbb{C}^*} & \text{fibré trivial sur $\mathbb{C}^*$ (fibre ponctuelle)}\\ \boxed{\mathcal{V}_1=\mathbb{D}} & \text{(germe de disques ouverts centrés en 0)}\\ \boxed{\mathcal{V}_1^*=\mathbb{D}^*} & \text{(germe de disques ouverts épointés en leur}\\ & \text{\quad centre 0)} \end{array}\right. \]
\note{la première ligne encadrée est prise entre deux traits ondulés, au-dessus et au-dessous\,; devant « germe », le mot « disque » est biffé.}
\[ \mathcal{V}_{1,0}\simeq\mathcal{V}_1,\qquad \mathcal{V}_{1,0}^*\simeq\mathcal{V}_1^* \]
\marginal{D\ill{} \ill{} notation \ill{} $D_{d,0}\xrightarrow{\sim}D_d$, \ill{} \ill{} que pour $d_*=(d_0,\dots,d_r)$, $d_r=0$, \ill{} \ill{} \ill{} \ill{} \ill{} $D_{d'_*}$}

En définitive, il y a « quatre » pièces de

\page{163}
construction » qui apparaissent
\[ (264)\quad \left\{\begin{array}{ll} \mathbb{C}^* & \text{(espaces analytiques « globaux »)}\\ \mathbb{P}\ \text{point} & \text{(\quad id\quad)}\\ \mathbb{D} & \text{germe de disques}\\ \mathbb{D}^* & \text{germe de disques épointés} \end{array}\right. \]
\note{une flèche courbe, à gauche de l'accolade, relie $\mathbb{P}$ à $\mathbb{C}^*$.}
dont deux « globaux », et deux « locaux », qui s'assemblent suivant le diagramme d'inclusions

\begin{tikzcd}
 & \mathbb{P} \arrow[d, hook] \\
\mathbb{D}^* \arrow[r, hook] \arrow[d, hook] & \mathbb{D} \arrow[d, dashed] \\
\mathbb{C}^* \arrow[r, dashed] & \mathbb{C}
\end{tikzcd}

\note{diagramme (265). Les deux flèches pointillées vers $\mathbb{C}$ et le sommet $\mathbb{C}$ sont entre crochets sur la page\,: le carré qu'elles ferment est celui que la phrase suivante décrit comme somme amalgamée.}

La variété analytique globale $\mathbb{C}$, avec son diviseur $\Theta$ localisé en $\mathbb{P}=\{0\}$, se reconstitue à partir de ce diagramme comme une somme amalgamée
\[ (266)\qquad \mathbb{C}\simeq\mathbb{C}^*\amalg_{\mathbb{D}^*}\mathbb{D}\,. \]

Si on veut comprendre la situation « de codimension $n$ » dans \struck{les} les espaces typiques les plus élémentaires, on prendra

\page{164}\note{p.~79 de l'auteur.}
\[ (267)\qquad X=\mathbb{C}^I,\qquad \Theta=\operatorname{div}\Big(\prod_{i\in I}z_i\Big) \]
\note{après la formule, un mot bref biffé, illisible.}
où $I$ est un ens. fini de cardinal $n$, et où $(z_i)_{i\in I}$ sont les fonctions coordonnées. On s'attend à ce que les « topos élémentaires » qui correspondent à cette situation soient les produits (indexés par $I$) de topos de l'un quelc. des quatre types envisagés dans (264) --- et que les morphismes canoniques entre ces topos doivent \struck{correspondre} \add{se déduire} des morphismes inclus dans le diagramme (265). Cela nous \struck{\ill{}} \uncertain{conduit}\,: considérons les partitions possibles de $I$ en quatre parties, correspondant aux quatre types en question. On écrira
\marginal{« partitions » au sens généralisé\,: il peut y avoir des « classes » vides.}
\[ (268)\quad \left\{\begin{array}{l} I=\underbrace{I'}_{\substack{\text{correspond aux facteurs}\\ \text{du type « global »}}}\amalg\underbrace{I''}_{\substack{\text{correspond aux facteurs}\\ \text{du type local}}}\\[18pt] I'=\underbrace{I'_0}_{\substack{\text{correspond aux facteurs } \mathbb{P}\\ \text{(de dim 0)}}}\amalg\underbrace{I'_1}_{\substack{\text{correspond aux facteurs } \mathbb{C}^*\\ \text{(de dim 1)}}}\\[18pt] I''=\underbrace{I''_0}_{\text{facteurs } \mathbb{D}}\amalg\underbrace{I''^*}_{\text{facteurs } \mathbb{D}^* \text{ (épointés)}} \end{array}\right. \]
\note{sous $I'_1$, le symbole lu $\mathbb{C}^*$ est surchargé et ressemble aussi à $\mathbb{D}^*$\,; le sens (facteurs « globaux » de dimension 1) appelle $\mathbb{C}^*$.}

\page{165}
Si
\[ d'_0=\operatorname{card} I'_0, \]
$d$ est la codimension (rel. à \struck{$X$} \add{$D_0$} $=\mathbb{C}^I$) du topos élémentaire envisagé, \struck{qui} \struck{ici dont} donc que
\[ d_1=\operatorname{card} I'_1 \]
\note{le symbole à gauche du signe égal est surchargé\,; on lit un $d$ d'indice 1, sans certitude.}
est la dimension de l'« âme » du topos, qui est par définition
\[ \text{\struck{$D^*_{I'_1}$}}=\underbrace{\prod_{i\in I'_1}\mathbb{C}^*}_{\mathbb{C}^{*I'_1}}\times\underbrace{\prod_{i\in I\setminus I'_1}\mathbb{P}}_{\mathbb{P}^{I'_0}} \]
\marginal{\struck{$\prod$ \ill{} \ill{} envisagé \ill{} $I$}}
\struck{On regarde $D_{I'_1}$ comme plongé dans}

\struck{On a, pour une partie $J$ de $I$,} \struck{$D_J=\mathbb{C}^{*J}\times\mathbb{P}^{I\setminus J}\subset X=\mathbb{C}^J$,} \struck{$D^*_J=\mathbb{C}^{*J}\times\mathbb{P}^{I\setminus J}$}\note{passage barré de quatre traits obliques.}

Pour une partie $J$ de $I$, on pose
\[ (269)\qquad \begin{array}{l} D_J=\mathbb{P}^J\times\mathbb{C}^{I\setminus J}\subset\mathbb{C}^I=X\\ \quad\cup\\ D^*_J=\mathbb{P}^J\times\mathbb{C}^{*\,I\setminus J} \end{array} \]
\marginal{sous-espace typique défini par l'annulation des coordonnées $z_j$, où $j\in J$, de codimension $d=\operatorname{card}J$.}
on aura donc
\[ (270)\qquad \left\{\begin{array}{l} D_d=\coprod_{J\subset\mathcal{P}(I),\ \operatorname{card}J=d} D_J\\ D^*_d=\coprod D^*_J \end{array}\right. \]

\page{166}\note{p.~80 de l'auteur.}
Avec cette notation, l'\emph{âme} du topos \add{élémentaire} envisagé est $D^*_{I\setminus I'_1}$, qu'on va considérer comme plongé dans $D^*_{I'_0}$\,:
\[ (271)\qquad \underset{(\text{« âme »})}{A}=D^*_{\underbrace{\scriptstyle I\setminus I'_1}_{I'_0\amalg I''}}\ \underset{\substack{\text{immersion}\\ \text{de codim. } I''}}{\hookrightarrow}\ D^*_{I'_0}=\underset{(\text{« réceptacle »})}{R} \]
Le \struck{premier} cas où on aurait $I''=I''_0$ correspond à celui d'un voisinage tubulaire de \struck{$D^*_A$ dans $D_{I'_0}$} $A$ dans $R$, soit $\mathcal{V}_{A,R}$), celui où $I''=I''^*$ correspond au cas de $\mathcal{V}^*_{A,R}$ (le $*$ bien sûr relatif ici à la situation induite sur $R$, où on y prend, sur $D_{I'_0}\supset R$, avec le diviseur induit par $\Theta_{d_0}$\,…). Le cas « général », correspond bien \uncertain{entendu} à la donnée d'une strate intermédiaire entre $A$ et $R$
\[ (272)\qquad \begin{array}{ccccc} A & \hookrightarrow & B & \hookrightarrow & R\\ \| & & \| & & \|\\ D^*_{I'_0\amalg I''} & \big| & D^*_{I'_0\amalg I''^*} & \big| & D^*_{I'_0}\\ & \scriptstyle\text{codim. card } I''_0=d''_0 & & \scriptstyle\text{codim.}=\operatorname{card}I''^*=d''^* & \end{array} \]

\page{167}
Je vais désigner par $\mathcal{V}^{B(*)}_{A,R}$ le topos correspondant, qui peut se décrire comme un produit fibré
\[ (273)\qquad \begin{array}{ccc} \mathcal{V}_{A,R} & \overset{\text{can}}{\hookrightarrow} & \mathcal{V}_{B,R}\\ \big\uparrow & & \big\uparrow\\ \mathcal{V}^{B(*)}_{A,R} & \hookrightarrow & \mathcal{V}^*_{B,R} \end{array} \]
Ce topos, dans la situation générale \struck{tout} précédente, dans les notations \struck{des} \uncertain{précédentes}\note{lecture d'ensemble de ces deux lignes incertaine\,; un mot biffé avant « notations ».} correspondant\,: \uncertain{données}
\[ (274)\qquad \begin{array}{ccccc} d''_* & \subset & d'_* & \subset & d_*\\ \| & & \| & & \|\\ \{d''_0>\dots>d''_{r''}\} & & \{d'_0>\dots>d'_{r'}\} & & \{d_0>\dots>d_r\} \end{array} \]
d'où
\[ (275)\qquad D_{d_*}\longrightarrow D_{d'_*}\longrightarrow D_{d''_*} \]
\struck{dans} on prendra le produit fibré \struck{$\mathcal{V}_{D_{d_*},D^{\ill{}}_{d''_*}}^{(*)}$} (qu'on notera par définition $\mathcal{V}^{d'_*}_{d_*,d''_*}$)\,:
\marginal{On va donc noter $\mathcal{V}^{d'_0}_{d_*,d''_0}$ ces topos.}
\note{la ligne inférieure du diagramme et la flèche verticale de droite sont celles de la page\,; les indices de la colonne de droite (prime ou non, $d$ ou $d_*$) sont lus avec peine.}

\begin{tikzcd}[column sep=small, row sep=small, nodes={font=\scriptsize}]
\mathcal{V}^*_{D_{d_*},D'_{d''_*}} \arrow[r, hook] & \mathcal{V}_{D'_{d_*},D_{d''_*}} \\
\mathcal{V}^{d'_*}_{d_*,d''_*} \arrow[u] \arrow[r] & \mathcal{V}^*_{D'_d,D''_d} \arrow[u, hook]
\end{tikzcd}

\note{diagramme (276) d'après le numéro de la marge (on lit « 270 », surchargé)\,; la flèche de droite monte vers $\mathcal{V}_{D'_{d_*},D_{d''_*}}$, et une seconde flèche courbe part de ce sommet vers une remarque en haut à droite\,: « ici les $d_*$ car\,\ill{} ».}
\marginal{[attention aux $*$ --- on ne \uncertain{s'en} occupe pas \uncertain{ici}\,: on veut pas en plus \uncertain{sur} $D_{d_*}$ \uncertain{l'action} $*$ --- l'âme doit être $D^*_{d_*}$\,!]}

\emph{NB} ce topos ne dépend en fait que des $d_*$, $d'_0$, $d''_0$, où $d_0\geqslant d'_0\geqslant d''_0\geqslant 0$\,?

\page{168}\note{p.~81 de l'auteur.}
Il me semble devenu plausible que ces topos, avec les morphismes canoniques entre eux qu'il faudrait expliciter, épuisent\,: expriment tous les topos \add{\uncertain{canoniques}} de \uncertain{stratification}, envisagés jusqu'à présent. \struck{Prenons} Prenons un cas-type, \add{non} \uncertain{élémentaire}, i.e. $\mathbb{C}^2$, \struck{\ill{}} il y a donc $4\times 4=16$ \struck{\ill{}} topos (qui s'ajoutent à un antérieur\,?)\note{la parenthèse et le mot « antérieur » sont lus avec peine\,; un trait relie « élémentaires » à « qui ».} de la forme\,:
\[ (277)\quad \left\{\begin{array}{ll} D_0^*=\mathbb{C}^*\times\mathbb{C}^* & (D_0=X=\mathbb{C}\times\mathbb{C})\\[2pt] \left.\begin{array}{l}\Delta_1^*=\mathbb{C}^*\times\mathbb{P}\\ \Delta_2^*=\mathbb{P}\times\mathbb{C}^*\end{array}\right\} & D_1^*=\Delta_1^*\amalg\Delta_2^*\quad (D_1=\Delta_1\amalg\Delta_2)\\ & \quad\text{où } \Delta_1=\mathbb{C}\times\mathbb{P},\ \Delta_2=\mathbb{P}\times\mathbb{C}\\[2pt] \mathbb{P}=\mathbb{P}\times\mathbb{P} & D_2^*=D_2=\mathbb{P}\\[6pt] \mathcal{V}_{\Delta_1^*,X}=\mathbb{C}^*\times\mathbb{D} & \mathcal{V}^*_{\Delta_1^*,X}=\mathbb{C}^*\times\mathbb{D}^*\\ \mathcal{V}_{\Delta_2^*,X}=\mathbb{D}\times\mathbb{C}^* & \mathcal{V}^*_{\Delta_2^*,X}=\mathbb{D}^*\times\mathbb{C}^*\\[2pt] \mathcal{V}_{\mathbb{P},\Delta_1}=\mathbb{D}\times\mathbb{P} & \mathcal{V}^*_{\mathbb{P},\Delta_1}=\mathbb{D}^*\times\mathbb{P}\\ \mathcal{V}_{\mathbb{P},\Delta_2}=\mathbb{P}\times\mathbb{D} & \mathcal{V}^*_{\mathbb{P},\Delta_2}=\mathbb{P}\times\mathbb{D}^*\\[6pt] \mathcal{V}_{\mathbb{P},X}=\mathbb{D}\times\mathbb{D}, & \\ \mathcal{V}^{\Delta_1}_{\mathbb{P},X}=\mathbb{D}\times\mathbb{D}^*, & \mathcal{V}^{\Delta_2}_{\mathbb{P},X}=\mathbb{D}^*\times\mathbb{D}\\ \mathcal{V}^*_{\mathbb{P},X}=\mathbb{D}^*\times\mathbb{D}^* & \end{array}\right. \]
\note{sur la page, une grande accolade marquée (277) embrasse le tout, et des accolades secondaires regroupent les lignes par types\,; un astérisque en marge marque le groupe des $D^*$. Les égalités $\mathcal{V}^*_{\Delta_2^*,X}$ et $\mathcal{V}^*_{\mathbb{P},\Delta_1}$ portent un astérisque surchargé.}
\marginal{Ces \ill{} écrites \ill{} \ill{} bien entendu \ill{} \ill{} \ill{} et symétriques sous les rangées \ill{} --- il faudrait \ill{} \ill{} \uncertain{écrire} \ill{}\,!}

\page{169}
Comment reconstituer p.\,ex. les \struck{ouverts} « globaux » \struck{partir} topos canoniques (cf. p.~74'\,\ill{}\note{p.~74' de l'auteur, hors de ce lot\,; le signe qui suit est illisible.}) qui sont des ouverts globaux de $X=\mathbb{C}^2$, savoir \struck{$\mathbb{C}^2\setminus\Delta_1$, \ill{}}
\[ (278)\qquad X=\mathbb{C}^2,\quad \underbrace{\mathbb{C}^2\setminus\Delta_1}_{U_1},\quad \underbrace{\mathbb{C}^2\setminus\Delta_2}_{U_2},\quad U=U_1\cup U_2=\mathbb{C}^2-\{\mathbb{P}\},\quad U_1\cap U_2=\mathbb{C}^*\times\mathbb{C}^* \]
Pour ces derniers, \struck{seuls} le dernier est déjà « élémentaire » --- c'est $D_0^*$. On aura
\[ \left.\begin{array}{l} U_1=\mathbb{C}\times\mathbb{C}^*\\ U_2=\mathbb{C}^*\times\mathbb{C}\end{array}\right\}\quad \mathbb{C}=\mathbb{C}^*\amalg_{\mathbb{D}^*}\mathbb{D} \]
\[ U=U_1\amalg_{U_1\cap U_2}U_2\qquad\text{où}\quad U_1\cap U_2=D_0^* \]
\[ X=\mathbb{C}\times\mathbb{C} \]
\struck{d'où il suffit d'exprimer}\quad $U_1$, $U_2$ seront donnés par les sommes amalgamées\note{« seront donnés » lu avec peine.}

\begin{tikzcd}[column sep=small, row sep=small, nodes={font=\scriptsize}]
 & \mathbb{D}^*\times\mathbb{C}^* \arrow[dl] \arrow[dr] & \\
\mathbb{C}^*\times\mathbb{C}^* \arrow[dr, dashed] & \text{loc.} & \mathbb{D}\times\mathbb{C}^* \arrow[dl, dashed] \\
 & U_1 & \\
 & \mathbb{C}^*\times\mathbb{D}^* \arrow[dl] \arrow[dr] & \\
\mathbb{C}^*\times\mathbb{C}^* \arrow[dr, dashed] & \text{loc.} & \mathbb{C}^*\times\mathbb{D} \arrow[dl, dashed] \\
 & U_2 &
\end{tikzcd}

\note{diagrammes (279), sous une accolade commune. Les flèches vers $U_1$ et $U_2$ sont tracées en pointillé sur la page, et chaque $U_i$ est entre parenthèses\,; la flèche $\mathbb{C}^*\times\mathbb{D}^*\to\mathbb{C}^*\times\mathbb{C}^*$ n'a pas de pointe nette.}

\page{170}\note{la page porte à nouveau le numéro 81 de l'auteur, comme la p.~168 du fonds.}
$U$ comme somme amalgamée

\begin{tikzcd}[column sep=small, row sep=small, nodes={font=\scriptsize}]
 & \mathbb{C}^*\times\mathbb{D}^* \arrow[dl] \arrow[dr] & & \mathbb{D}^*\times\mathbb{C}^* \arrow[dl] \arrow[dr] & \\
\mathbb{C}^*\times\mathbb{D} \arrow[drr, dashed] & & \mathbb{C}^*\times\mathbb{C}^* & & \mathbb{D}\times\mathbb{C}^* \arrow[dll, dashed] \\
 & & U=\mathbb{C}^2\setminus\{\mathbb{P}\} & &
\end{tikzcd}

\note{diagramme (280)\,; les deux flèches vers $U$ sont en pointillé et $U$ est entre parenthèses, comme en (279).}

$\mathbb{C}\times\mathbb{C}$ comme somme amalgamée du diagramme suivant (de type produit de $\longleftarrow\!\cdot\!\longrightarrow$ par lui-même)
\struck{$\mathbb{C}^*\times\mathbb{D}^*\leftarrow\mathbb{D}^*\times\mathbb{D}^*\to\mathbb{D}^*\times\mathbb{D}^*$, $\mathbb{D}^*\times\mathbb{D}^*$}\note{première ébauche du diagramme, raturée de traits en zigzag\,; la lecture des sommets est incertaine.}

\begin{tikzcd}[column sep=small, row sep=small, nodes={font=\scriptsize}]
\mathbb{C}^*\times\mathbb{C}^* & \mathbb{D}^*\times\mathbb{C}^* \arrow[l] \arrow[r] & \mathbb{D}\times\mathbb{C}^* \\
\mathbb{C}^*\times\mathbb{D}^* \arrow[u, no head] \arrow[d] & \mathbb{D}^*\times\mathbb{D}^* \arrow[l] \arrow[r] \arrow[u] \arrow[d] & \mathbb{D}\times\mathbb{D}^* \arrow[u] \arrow[d, no head] \\
\mathbb{C}^*\times\mathbb{D} & \mathbb{D}^*\times\mathbb{D} \arrow[l] \arrow[r] & \mathbb{D}\times\mathbb{D}
\end{tikzcd}

\note{diagramme (281). Sur la page, $\mathbb{D}\times\mathbb{D}^*=\mathcal{V}^{\Delta_2}_{\mathbb{P},X}$ et, sous $\mathbb{D}^*\times\mathbb{D}$, $\|\ \mathcal{V}^{\Delta_1}_{\mathbb{P},X}$. Deux traits verticaux (colonne de gauche vers le haut, colonne de droite vers le bas) n'ont pas de pointe visible.}
\marginal{Somme amalgamée $=\mathbb{C}^2$}

\page{171}
C'est dans ces deux diagrammes qu'interviennent les « non-locaux » topos élémentaires qui n'interviennent dans « l'introduction », savoir le fameux $\mathcal{V}^{\Delta_1}_{\mathbb{P},X}$ et $\mathcal{V}^{\Delta_2}_{\mathbb{P},X}$.\note{la lecture « qui n'interviennent dans » est douteuse\,; la phrase n'a peut-être pas de négation complète.}

Je voudrais maintenant expliciter le voisinage des topos $\mathcal{V}^{d'_0}_{d_*,d''_0}$ pour $d'_0,d''_0\in d_*$, $d'_0\geqslant d''_0$ \struck{\ill{}} \struck{($d''_*\subset d'_*\subset d_*$)} variable, et celle des $\pi_1$ associés, et expliciter comment, à l'aide de ces topos sur \uncertain{groupoïdes}, on peut reconstituer tous les « topos canoniques » associés\,: la stratification envisagée --- ou leurs groupoïdes fondamentaux.

Avant \struck{de \ill{}} d'examiner \add{\uncertain{spécialement}} le voisinage des $\mathcal{V}^{d'_0}_{d_*,d''_0}$, je voudrais \struck{\ill{}} regarder un peu la structure des « espaces » \emph{homotopiques} précédents, le fibré normal du $D_{d_*}$ \struck{dans $D_{d''_0}$, \ill{} $\check{N}_{d_*,d'_*}$} \struck{$(282)$ $\check{N}_{d_*,d''_0}\overset{\text{déf}}{=}\check{N}_{D_{d_*},D_{d''_0}}$ fibré normal}\note{les deux dernières lignes, avec le numéro (282), sont barrées d'un trait horizontal et d'un trait oblique\,; « déf » est écrit au-dessus du signe égal.}

\emph{NB}

\page{172}\note{p.~82 de l'auteur. La page commence par la fin de la phrase ouverte au bas de la p.~171.}
pour l'immersion locale
\[ D_{d_*}\longrightarrow D_0 \]
soit $\check{N}_{d_*}$, il y a dans $\check{N}_{d_*}$ une filtration canonique
\[ (282)\qquad \begin{array}{ccccccccccc} \check{N}_{d_*} & \overset{\text{déf}}{=} & \check{N}_{d_*,d_0} & \supset & \check{N}_{d_*,d_1} & \supset & \check{N}_{d_*,d_2} & \supset\cdots\supset & \check{N}_{d_*,d_r} & \supset & \check{N}_{d_*,0}\\ & & \|\,{\scriptstyle\text{déf}} & & & & & & & & \|\\ & & \check{N}_{D_{d_*},D_{d_0}} & & & & & & & & \{0\}_{D_{d_*}}\\ & & \text{rg } d_0 & & \text{rg } d_1 & & \text{rg } d_2 & & \text{rg } d_r & & \end{array} \]
\marginal{\emph{NB} $d_*=\{d_0>d_1>\dots>d_r>0\}$}
où $\check{N}_{d_*,d_i}$ est défini comme le \struck{fibré conormal \ill{}} \struck{\ill{} la suite exacte de fibrés} \struck{conormaux associés à la suite d'immersions} \struck{$D_{d_*}\to D_{\{d_i\}}\to D_0$}\note{trois lignes barrées d'un trait ondulé et de traits horizontaux\,; la lecture des mots biffés est incertaine.}
\[ (283)\qquad \check{N}_{d_*,d_i}=\check{N}_{D_{d_*},D_{d_i}}, \]
\struck{$(\check{N}_{d_i\to D_0}|D_{d_*})\to\check{N}$}\quad dans la suite exacte \struck{$\|$ déf} canonique \struck{\ill{}} associée à l'immersion locale $D_{d_*}\to D_{d_i}$\,:
\[ (284)\qquad 0\longrightarrow\check{N}_{d_*,d_i}\longrightarrow\check{N}_{d_*}\longrightarrow(\check{N}_{d_i}\,|\,D_{d_*})\longrightarrow 0\,; \]
plus gén., si $d'_*\subset d_*$ et qu'on prenne $d_i=d'_0$, on déduit de (284) la forme
\[ (285)\qquad 0\longrightarrow\underset{\substack{\wr\\ \check{N}_{d_*,d'_0}}}{\check{N}_{D_{d_*},D_{d'_*}}}\longrightarrow\underset{\substack{\|\,\text{déf}\\ \check{N}_{d_*}}}{\check{N}_{D_{d_*},D_0}}\longrightarrow(\check{N}_{D_{d'_*}}\,|\,D_{d_*})\longrightarrow 0 \]
\note{la page qualifie ces fibrés de « conormaux » dans les lignes biffées, mais le texte non biffé de la p.~171 parle de « fibré normal »\,; la notation $\check{N}$ est la sienne. Dans (282), $\check{N}_{d_*,0}$ est surchargé (on lit peut-être « $\check{N}_{d_*,0}$ » sur un autre indice).}

\page{173}
Le quotient canonique le plus général associé à la filtration (282) est le fibré vectoriel
\[ (286)\qquad \check{N}_{d_*;(d'_0,d''_0)}\overset{\text{déf}}{=}\check{N}_{d_*,d''_0}\big/\check{N}_{d_*,d'_0} \]
où on suppose
\[ d'_0,\ d''_0\in d_*,\qquad d''_0\leqslant d'_0\ (\leqslant d_0), \]
on a donc une suite exacte canonique
\[ (287)\qquad 0\longrightarrow\check{N}_{d_*,d'_0}\longrightarrow\check{N}_{d_*,d''_0}\longrightarrow\check{N}_{d_*;(d'_0,d''_0)}\longrightarrow 0\,. \]
\note{l'inégalité $d''_0\leqslant d'_0$ s'accorde mal avec l'inclusion $\check{N}_{d_*,d'_0}\subset\check{N}_{d_*,d''_0}$ que demande (287) au vu de la filtration décroissante (282)\,; la page porte bien ce qui est transcrit, sous réserve du sens des primes, surchargés par endroits.}
Cette suite exacte peut s'interpréter ainsi\,: on a
\[ \begin{array}{ccccc} d''_* & \subset & d'_* & \subset & d_*\\ \| & & \| & & \\ \{d''_0>\dots>d''_{r''}\} & & \{d'_0>\dots>d'_{r'}\} & & \end{array} \]
donc la suite exacte
\[ (288)\qquad \begin{array}{ccccccccc} 0 & \to & \check{N}_{D_{d_*},D_{d'_*}} & \to & \check{N}_{D_{d_*},D_{d''_*}} & \to & (\check{N}_{D_{d'_*},D_{d''_*}}\,|\,D_{d_*}) & \to & 0\\ & & \wr & & \wr & & \wr & & \\ & & \check{N}_{d_*,d'_0} & & \check{N}_{d_*,d''_0} & & \check{N}_{d_*;(d'_0,d''_0)} & & \end{array} \]
Donc, dans ce cas, un fibré vectoriel $\check{N}/$\,(celui-là le plus général\,\ill{}) \uncertain{est} \uncertain{dans} $\check{N}_{d_*;(d'_0,d''_0)}$\note{fin de page incertaine\,; un trait oblique et une parenthèse en marge droite semblent renvoyer à la suite, hors de ce lot ou à la page suivante.}

\page{174}\note{p.~83 de l'auteur.}
\ill{} on a loc. une décomposition en produits de fibrés en droites, d'où (en prenant les fibrés principaux associés) un ouvert
\[ (289)\qquad N^*\subset N,\qquad\text{en particulier}\quad \left\{\begin{array}{l} N^*_{d_*}\subset N_{d_*}\\ N^*_{d_*,d_i}\subset N_{d_*,d_i}\\ N^*_{d_*;(d'_0,d''_0)}\subset N_{d_*;(d'_0,d''_0)} \end{array}\right. \]
qui est d'ailleurs un \emph{fibré} torseur sous tores convenables, que je me dispense d'expliciter.
\note{ici et dans la suite de la page, le chevron sur $N$ est tantôt présent, tantôt absent\,; la transcription suit la page.}

Considérons alors
\[ (290)\qquad \check{N}^{d'_0}_{d_*,d''_0}=\text{image inverse de }\check{N}^*_{d_*;(d'_0,d''_0)}\text{ dans }\check{N}_{d_*,d''_0}, \]
par l'homomorphisme
\[ \check{N}_{d_*,d''_0}\longrightarrow\check{N}_{d_*;(d'_0,d''_0)}\quad\text{de (287)}. \]
On a dans ce cas général\,:
\[ (291)\qquad \check{N}^{d'_0}_{d_*,d''_0}\ \text{est fibré sur}\ \check{N}^*=\check{N}^*_{d_*;(d'_0,d''_0)}, \]
et est un torseur sous le fibré vectoriel sur $\check{N}^*$, image inverse de $\check{N}_{d_*,d_0}$ par $\check{N}^*\to D_{d_*}$.\note{l'indice de $\check{N}_{d_*,d_0}$ est peut-être $d'_0$\,; celui de $D_{d_*}$ porte un signe surchargé.}

\page{175}
Je dis qu'on a alors une \struck{équivalence} « homéomorphisme » (heuristique\,!)
\[ (292)\qquad \check{N}^{d'_0}_{d_*,d''_0}\,|\,D^*_{d_*}\ \simeq\ \mathcal{V}^{d'_0}_{d_*,d''_0} \]
qui induit, grâce à l'homomorphisme
\[ (293)\qquad \check{N}^{d'_0}_{d_*,d''_0}\,|\,D^*_{d_*}\longrightarrow\check{N}^*_{d_*;(d'_0,d''_0)}\,|\,D^*_{d_*} \]
qui est une \emph{équivalence d'homotopie}, une \emph{équivalence d'homotopie}
\[ (294)\qquad \boxed{\ \mathcal{V}^{d'_0}_{d_*,d''_0}\ \simeq\ \check{N}^*_{d_*;d'_0,d''_0}\,|\,D^*_{d_*}\quad\text{équivalence d'homotopie}\ } \]
qui a le sens d'un « voisinage » de lisse\note{« voisinage de lisse » lu avec peine.} par correspondance à une fléchée des \uncertain{lacets} \uncertain{canoniques} dans \uncertain{une} \uncertain{seule} où les \uncertain{lacets} \ill{}, \ill{} (par la \uncertain{\ill{}}) d'avoir une \uncertain{relation} \uncertain{précise} \emph{homotopique}. Retenons cependant la
\note{la phrase s'interrompt au bas de la page\,; la lecture des quatre dernières lignes est très incertaine.}

\page{176}\note{p.~84 de l'auteur.}
\emph{Scholie important.} Les « pièces de construction » dans la dissection d'un espace analytique [lisse], muni d'un diviseur à croisements normaux, sont, à homotopie près, les espaces obtenus ainsi, fibrés sur les $D^*_{d_*}=D^*_{(d_0,d_1,\dots,d_r)}$ (qui sont analytiques lisses, \struck{\ill{}} complémentaires de partie finie \uncertain{finie} \struck{\ill{}} de $\mathcal{N}$)\note{la parenthèse est lue avec peine\,; « $\mathcal{N}$ » est peut-être un $\Theta$ ou un autre symbole.} ainsi\,: on regarde \struck{\ill{}} $\check{N}_{d_*}|D^*_{d_*}$ sur $D^*_{d_*}$ rel. \ill{}, (\ill{} $D_{d_*}\to D_{d_0}$) en fibré \add{vectoriel} \struck{\ill{}} muni d'une filtration du type $d_*$, on prend un quotient général $\check{N}=\check{N}_{d_*;(d'_0,d''_0)}$ de cette filtration, \struck{\ill{}} (contractile\,: $D^*_{d_*}$), et l'ouvert $\check{N}^*$ --- qui est donc un fibré loc. trivial sur $D^*_{d_*}$, à fibres $\mathbb{C}^{*\,d'_0-d''_0}$ --- c'est bien\,! En \uncertain{fatigant} un peu, on doit
\marginal{\ill{} compris le cas \ill{} \ill{} \ill{} associés \ill{} \ill{} \ill{} sur $D^*_0$ \ill{} \ill{} fibré \ill{} le fibré \ill{} $D_0$ \ill{} \ill{} trivial \ill{} \ill{} \ill{} \ill{} l'\ill{} \ill{} \ill{} \ill{} $D_0$ \ill{} $\Theta$ \ill{}}
\note{la note marginale est écrite en diagonale dans l'angle gauche de la page, en lignes serrées\,; on n'en reconnaît que quelques mots.}

\page{177}
écrire en tous cas\,: définir une équivalence \emph{canonique} sur les groupoïdes fondamentaux des $\check{N}^*_{d_*,(d'_0,d''_0)}|D^*_{d_*}$ d'une part, les $\mathcal{V}^{d'_0}_{d_*,d''_0}$ d'autre part
\[ (295)\qquad \boxed{\ \Pi_1\big(\underbrace{\check{N}^*_{d_*;(d'_0,d''_0)}\,|\,D^*_{d_*}}_{\overset{\text{déf}}{=}\ \check{N}^{**}_{d_*;(d'_0,d''_0)}}\big)\xrightarrow{\ \approx\ }\Pi_1\mathcal{V}^{d'_0}_{d_*,d''_0}\ } \]
\struck{Accessoirement}\note{mot biffé, lecture incertaine.} Actuellement on \uncertain{prend} la connaissance du $\Pi_1\mathcal{V}^{d'_0}_{d_*,d''_0}$ (\uncertain{modulo} équivalence) \uncertain{revient} donc à celle des $\Pi_1$ extensions de $\check{N}^{**}$ \uncertain{fibré} au-dessus des affines\note{« affines » lu avec peine.} composantes connexes de $D^*_{d_*}$ --- et on sait, par la suite exacte d'homotopie, des extensions des $\Pi_1$ des comp. conn. de $D^*_{d_*}$, par des quotients en $\mathbb{Z}^r$ (par l'image du $\Pi_2$ desdites comp. connexes\,…)

\page{178}\note{p.~85 de l'auteur.}
\emph{Remarque.} Lorsque les $D^*_{d_*}$ sont tous des $K(\pi,1)$ (ce qui est le cas dans la situation de Teichmüller\,!) alors la connaissance des systèmes des $\Pi_1\mathcal{V}^{d'_0}_{d_*,d''_0}$ précédents (avec les morphismes canoniques entre ces groupoïdes) doit permettre de reconstituer entièrement le \struck{type d'homotopie} \add{groupoïde} de $X$ et des \struck{topos} \add{divers} topos canoniques associés à $X$, mais aussi le topos quasi-galoisien \uncertain{formé} par les faisceaux sur les topos qui sont des quotients de topos « élémentaires » constitutifs, i.e. loc. constants --- et par le biais le type d'homotopie des topos envisagés --- y compris $X$\,…
\marginal{(par \uncertain{définition}, \ill{} \ill{} \ill{} \ill{} \ill{} \ill{} $K(\pi,1)$, \uncertain{faisceaux} \ill{} \ill{} \ill{}\,…)}
\note{la note marginale, écrite en diagonale dans la marge gauche et reliée par une flèche au mot « constitutifs », n'est lisible que par fragments.}

\emph{Remarque.} Ce n'est pas la compréhension du topos \uncertain{élémentaire}, pas un \uncertain{revendu}\note{fin de page\,; la phrase continue hors de la page, la lecture de la dernière ligne est très incertaine.}

\page{179}
compte --- le cas le plus simple de $X\simeq\mathbb{P}^1_{\mathbb{C}}\approx S^2$, avec un diviseur lisse réduit à un point, donc $X^*\simeq\mathbb{C}$ (qui est bien un $K(\pi,1)$ avec $\pi=\{1\}$\,!), est le contre-exemple type, \struck{\ill{}} \struck{car} on ne reconstitue pas le type d'homotopie de $X=S^2$ \struck{à partir} \add{décrit} via les systèmes « inductifs » des topos galoisiens associé\supplied{s}\,:
\note{« décrit » est écrit au-dessus d'un mot biffé, peut-être « exhibé »\,; un autre mot biffé, au-dessus de la ligne, se lit « homotopie » sans certitude.}

\begin{tikzcd}[column sep=small, row sep=small]
 & \mathbb{D}^*_\infty \arrow[dl] \arrow[dr] & \\
\mathbb{D}_\infty & & {\mathbb{C}=\mathbb{P}^1_{\mathbb{C}}\setminus\{\infty\}=X^*}
\end{tikzcd}

\note{sous $\mathbb{C}$, le sommet de droite porte une rature épaisse\,; on lit $\mathbb{C}$ au-dessous, avant l'égalité.}
i.e. au diagr.\ de groupes \struck{\ill{}}

\begin{tikzcd}[column sep=small, row sep=small]
 & \mathbb{Z} \arrow[dl] \arrow[dr] & \\
\{1\} & & \{1\}
\end{tikzcd}

\struck{qui se réalise} \struck{\ill{}} le topos global « galoisien » associé (faisceaux sur $S^2$ qui sont loc. constants, donc constants sur $S^2\setminus\{\infty\}=\mathbb{C}$) s'identifie à $\hat{C}$, où $C$ est la catégorie $\to\cdot$ qui a un seul point, donc la cohomologie de $\hat{C}$ est nulle en dim $>0$, alors que celle de $S^2$ ne l'est pas. Ceci souligne la question générale de comprendre le type d'homotopie d'un topos obtenu par « recollement » d'un ouvert à son fermé complémentaire\,…
\note{« recollement d'un ouvert à son fermé complémentaire » est lu avec peine\,; la flèche $\to\cdot$ désigne vraisemblablement la catégorie à deux objets et une flèche, mais la page ne précise pas.}

\page{180}\note{p.~86 de l'auteur. Nouveau paragraphe, numéroté 13 par l'auteur\,; il se poursuit au-delà de ce lot.}
\section{13. Digressions sur stratifications, locales et globales}

\emph{Une stratification globale} d'un espace topologique $X$ (pour fixer les idées\,!) est une partition de $X$ en « strates » $(X_i)_{i\in I}$\note{« strates » est écrit au-dessus d'un mot biffé, peut-être « cellules ».} ayant les propriétés\,:

a) les $X_i$ sont loc. fermés --- la famille est loc. finie\note{« $X_i$ » et « est loc. finie » sont ajoutés sous la ligne\,; la construction exacte reste douteuse.}

b) $\forall i\in I$, $X_i$ est une réunion de \struck{\ill{}} \add{strates} loc. connexes et \uncertain{connexes}\note{« loc. connexes et » est ajouté au-dessous, « connexes » encadré d'un trait\,; au-dessus, « elles mêmes connexes » lu avec peine.} (si \uncertain{tel}, on suppose\,…)

c) Les strates sont des façons plus fortes que b), que \struck{\ill{}} \add{les} composantes connexes \struck{d'une} $X_{i,\alpha}$ d'un $X_i$, et $\forall X_j$, $\overline{X_{i,\alpha}}\cap X_j$ est une réunion de comp. connexes de $X_j$\,; donc on remplaçant la partition par les $X_i$ par celle par les $X_{i,\alpha}$, on satisfait\,: a) b) c)\,…)

d) \emph{Les strates} $\forall X_i$, et tout $x\in\overline{X_i}$, le germe d'\uncertain{adhérence} des $X_i$ en un voisinage de $x$ est connexe (ceci exclut la stratification associée à la courbe à croisement normaux
\drawing{une boucle\,: une courbe qui se recoupe en un point double.}
ou celle-ci\,:
\drawing{un disque hachuré, deux points marqués sur son bord.}
(une strate de dim. $0$, une de dim. $1$, \struck{\ill{}} une de dim. $2$)\note{parenthèse écrite en petits caractères sous le dessin\,; un mot biffé avant « une de dim. 2 ».}). Plus précisément, on veut ceci\,:

e) $\forall X_i$ la structure locale autour de $X_i$ \add{de l'adh.} des $\overline{X_j}$ tels que $\overline{X_j}\supset X_i$ est \uncertain{loc.}
\note{la page s'arrête au milieu de la condition e)\,; la suite est hors de ce lot.}

\end{document}
